problem:  
Let \((M^{2n+1}, T^{1,0}M, \theta)\) be a compact, strictly pseudoconvex CR manifold with nonnegative CR Yamabe invariant \(Y(M)\ge 0\).  
Consider the CR Yamabe flow  
\[
\frac{\partial \theta_t}{\partial t}=-(R_{\theta_t}-\bar{R}_{\theta_t})\theta_t,\qquad \theta|_{t=0}=\theta_0,
\]
where \(R_{\theta_t}\) denotes the Webster scalar curvature and \(\bar{R}_{\theta_t}\) its average over \(M\).  
Assume \(R_{\theta_0}\) changes sign, and the flow develops a finite-time singularity at \(T<\infty\).

**Problem:**  
Suppose that for some sequence \(t_i\to T\) and points \(p_i\in M\), there exist parabolic scaling factors \(\lambda_i>0\) such that the rescaled pseudo-Hermitian structures  
\[
\tilde{\theta}_i = \lambda_i^2 \theta_{t_i},
\]
satisfy \(\max_M R_{\tilde{\theta}_i}=1\) and that the pointed manifolds \((M,p_i,\tilde{\theta}_i)\) subconverge (after passing to a subsequence and modulo CR diffeomorphisms) in the pointed \(C^\infty_\text{loc}\)-topology to a complete pseudo-Hermitian manifold \((M_\infty,p_\infty,\theta_\infty)\).  

**Classify** all possible limit geometries \((M_\infty,\theta_\infty)\) that can arise in this situation.  
In particular, does every such limit have Webster scalar curvature \(R_{\theta_\infty}\equiv 1\) and CR structure locally equivalent to either:  
(a) the standard CR sphere \((S^{2n+1},\theta_{\mathrm{std}})\), or  
(b) the Heisenberg group \((\mathbb{H}^{2n+1},\theta_{\text{H}})\)?  

Alternatively, can there exist nontrivial CR Yamabe solitons as other possible singularity models?  

Why is it a "good" problem:  
This question precisely isolates a deep analytic and geometric frontier in the study of the CR Yamabe flow. It formulates, under minimal assumptions, the **classification problem for possible singularity models**—an analogue of Hamilton’s classification of Ricci flow blow-ups. The hypotheses are standard, and the compactness/convergence condition is clearly stated, separating existence of blow-up limits from their classification. The potential discovery or exclusion of new CR Yamabe solitons would illuminate how curvature, conformal invariance, and CR automorphism symmetry control singularity formation. This bridges geometric analysis, PDE theory, and CR geometry in the same spirit as Ricci flow theory connects geometry to parabolic analysis. The problem’s simplicity in expression and its capacity to open a new theory of CR singularities make it a “good” research problem in the sense of being both foundational and forward-looking.